%! AP Statistics — Unit 3, Lesson 3.5 — Group Activity
\documentclass[10pt]{article}
\usepackage{apstats-article}
\usepackage{apstats-boxes}

\begin{document}

\pageheader{Unit 3, Lesson 3.5}{Group Activity: Designing Experiments}
\namepartnerperiod

\vspace{0.1in}

\begin{scenariobox}[Context]{sky}
\small
Your group will analyze and design experiments. For each scenario, identify the
components, evaluate whether the four principles are met, and improve the design
where needed. Then design your own experiment from scratch.
\end{scenariobox}

\vspace{0.1in}

\textbf{Scenario A:} A researcher wants to know if a new supplement improves memory.
She recruits 30 volunteers. Those who ask to take the supplement get it; the rest get
nothing. After 4 weeks, all participants take a memory test.

\begin{practicebox}
\small
\begin{enumerate}[leftmargin=1.5em, itemsep=7pt]
  \item Experimental units: \blank{8cm}
  \item Factor: \blank{8cm}
  \item Treatments: \blank{8cm}
  \item Response variable: \blank{8cm}
  \item Which of the four principles is \emph{violated}? \blank{5cm} Why?\\
        \writeline\\[1pt]\writeline
  \item What design flaw allows a confounding variable? Describe it.\\
        \writeline\\[1pt]\writeline
  \item Redesign: How would you fix this experiment?\\
        \writeline\\[1pt]\writeline
\end{enumerate}
\end{practicebox}

\vspace{0.1in}

\textbf{Scenario B:} Researchers test two brands of headache medication. They have
40 patients with chronic headaches. Half are randomly assigned to Brand A;
half to Brand B. Neither patients nor doctors know which brand a patient is taking.
After 2 weeks, doctors rate each patient's improvement.

\begin{practicebox}
\small
\begin{enumerate}[leftmargin=1.5em, itemsep=7pt]
  \item Design type (CRD, block, matched pairs): \blank{6cm}
  \item What type of blinding is used? \blank{6cm}
  \item Is there a control group? \blank{3cm} Should there be? Explain.\\
        \writeline\\[1pt]\writeline
  \item Patients vary in age (some are 20--30 years old; some are 60--70). Would
        blocking by age group improve this design? Explain.\\
        \writeline\\[1pt]\writeline
\end{enumerate}
\end{practicebox}

\vspace{0.1in}

\textbf{Part 3: Design your own experiment.}

\begin{practicebox}
\small
Research question: \textit{Does listening to classical music while studying improve
quiz scores for AP Statistics students?}

\vspace{6pt}
Using 30 AP Statistics students, design a well-controlled experiment. Address:
\begin{enumerate}[leftmargin=1.5em, itemsep=6pt]
  \item Treatments (be specific): \blank{9cm}
  \item How you will randomly assign students: \blank{9cm}
  \item Control of confounding variables (list 2 and how you address each):\\
        1. \blank{9cm}\\[4pt]
        2. \blank{9cm}
  \item Response variable and how it will be measured: \blank{9cm}
  \item Will you use blinding? What type is feasible? \blank{9cm}
  \item Draw a diagram of your design in the space below.
  \vspace{1.6in}
\end{enumerate}
\end{practicebox}

\end{document}
